\bookchapter{附录F：折叠屏选购与使用专题}

\chapterintro{折叠屏已经从“未来科技”变成“现在能买”的成熟品类。但它和直板机在形态、使用习惯、保养方式上差异很大，值不值得买、怎么选、怎么用，这一篇帮你理清楚。}

\booksection{一、折叠屏现在发展到什么程度了？}

先说结论：\textbf{折叠屏已经不是“尝鲜玩具”，但也不是“人人该买”的主流形态}。它介于小众极客产品和大众消费品之间，适合特定人群。

目前主流分两类：

\begin{itemize}
  \item \textbf{横向折叠（书本式）}：合上是普通手机大小，展开是接近方形的小平板。代表：华为Mate X系列、三星Galaxy Z Fold系列、小米MIX Fold系列、vivo X Fold系列。主打"一部手机=手机+平板"。
  \item \textbf{纵向折叠（翻盖式）}：合上像粉饼盒，展开是正常尺寸手机。代表：华为Pocket系列、三星Galaxy Z Flip系列、OPPO Find N Flip、小米MIX Flip。主打"小巧便携+外屏速览"。
\end{itemize}

2024-2026年的关键进步：铰链寿命普遍标称50万次折叠（每天开合100次能用13年）、UTG超薄玻璃取代早期塑料膜、折痕肉眼可见但手感已很轻微、价格从万元档下探到4000-5000元档。

\begin{tcbinfo}{一句话判断}
如果你现在用直板机完全够用、没有“屏幕不够大”的痛点，就别买折叠屏。它是“锦上添花”的设备，不是“雪中送炭”的设备。

\end{tcbinfo}

\booksection{二、什么人适合买折叠屏？}

根据真实用户反馈，这几类人买折叠屏“真香”概率高：

\begin{enumerate}
  \item \textbf{重度阅读/文档处理者}：经常在手机上看PDF、批注文档、读长文。展开后接近7.5英寸的内屏，看A4 PDF不用缩放，体验质变。
  \item \textbf{经常出差/开会的商务用户}：不想带平板又嫌手机屏小。横向折叠屏展开看PPT、签合同、双屏对照邮件，比直板机效率高一截。
  \item \textbf{喜欢"悬停态"创作的人}：折叠到90度放在桌上，上半屏取景下半屏按快门（自拍/延时摄影）、上半屏视频下半屏记笔记、上半屏直播下半屏看评论——这是直板机做不到的形态。
  \item \textbf{追求便携又想要大屏的女性用户}：纵向折叠屏合上能塞进小包，展开是正常手机，外屏还能当化妆镜、看通知。
\end{enumerate}

\begin{tcbwarn}{这几类人别买}
\begin{itemize}
  \item 手机只用来刷短视频、聊微信、打电话的——折叠屏对你毫无增益，反而更重更厚。
  \item 手小或手腕有腱鞘炎的——横向折叠屏展开后普遍230g+，单手握持累。
  \item 预算紧张、一部手机要用4年以上的——折叠屏二手保值率低于直板机，维修成本高。
\end{itemize}

\end{tcbwarn}

\booksection{三、选购要点：看这6个指标}

别被“折叠”两个字晃花了眼，选购时重点看：

\begin{itemize}
  \item \textbf{铰链类型}：水滴铰链（折痕浅、寿命长，目前主流）优于U型铰链（折痕深、早期方案）。问销售或查评测确认是哪种。
  \item \textbf{内屏材质}：必须是UTG超薄玻璃，不要早期那种塑料膜。UTG手感接近玻璃，塑料膜指甲划过会有凹陷感且易留痕。
  \item \textbf{折痕表现}：去线下店亲手展开看。2026年的合格水平是：亮屏时正面几乎看不出，斜光下能看出一条浅沟，手指划过有轻微凹陷但不割手。
  \item \textbf{外屏尺寸}：横向折叠屏的外屏越大越好（最好6.4英寸以上），合上能当正常手机用；纵向折叠屏的外屏越大越实用（4英寸以上能回消息、看导航）。
  \item \textbf{重量与厚度}：横向折叠展开后230-260g是常态，低于230g算优秀；合上厚度12-14mm，低于12mm算薄。纵向折叠合上后能轻松放进衬衫口袋才算合格。
  \item \textbf{售后与保修}：重点看内屏碎裂的维修价格。有些品牌首年免费换一次内屏，有些自费要2000-4000元。这是折叠屏最大的隐性成本。
\end{itemize}

\begin{tcbpractice}{线下体验清单}
\begin{enumerate}
  \item 亲手开合10次，感受铰链阻尼是否顺滑、有没有异响。
  \item 展开后斜对着灯光看折痕，再正面亮屏看。
  \item 单手握持展开状态30秒，感受重量是否能接受。
  \item 试用悬停态：折90度放桌上，试试拍照和看视频。
  \item 问清楚内屏单独维修价格和保修政策。
\end{enumerate}

\end{tcbpractice}

\booksection{四、折叠屏独有的使用技巧}

买回来之后，这几个玩法是直板机体验不到的：

\booksubsection{分屏与多任务}

横向折叠屏展开后，大部分系统支持上下/左右分屏。典型场景：上半屏看网课视频，下半屏记笔记；左边微信右边文档；上面地图导航下面音乐。比直板机分屏实用得多，因为屏幕够大，每半边都不挤。

\booksubsection{悬停态（90度折叠）}

这是折叠屏的“杀手锏”用法：

\begin{itemize}
  \item \textbf{自拍/合影}：放桌上折90度，上半屏取景，下半屏按快门，不用找人帮忙拍。
  \item \textbf{延时摄影/夜景}：同样放桌上稳定拍摄，省三脚架。
  \item \textbf{视频通话}：解放双手，对方看你角度也更自然。
  \item \textbf{看视频}：放桌上当小电视，飞机/高铁小桌板上特别实用。
  \item \textbf{外屏预览+内屏拍摄}：用后置摄像头拍自己时，外屏能实时显示画面，构图更准。
\end{itemize}

\booksubsection{外屏速览}

纵向折叠屏合上后，外屏可以：看时间/通知/天气、回微信（语音转文字或快捷回复）、刷公交卡、看导航、当化妆镜（调用前置摄像头）。不用每次都展开，省电又省事。

\booksubsection{电脑模式}

部分横向折叠屏（华为、三星）展开后可切换“电脑模式”，外接显示器键鼠后像台简易PC，临时处理文档、回邮件够用。不是主力生产力工具，但应急很香。

\booksection{五、保养要点：让折叠屏活得久一点}

折叠屏最娇贵的就是内屏和铰链，这几条习惯能显著延长寿命：

\begin{itemize}
  \item \textbf{别在低温下强行开合}：0℃以下UTG玻璃变脆，开合前先揣怀里暖几分钟。这是官方说明书都会写但没人看的重点。
  \item \textbf{避免指甲/硬物接触内屏}：内屏虽有保护层但比外屏软，用指腹操作，别用指甲戳。
  \item \textbf{合上时确认屏内无异物}：一粒沙子夹在中间合上，可能直接压坏内屏。合上前扫一眼。
  \item \textbf{不要贴普通钢化膜}：内屏只能用官方专用膜或UTG保护膜，普通钢化膜会影响折叠和触控。建议去官方售后贴膜。
  \item \textbf{铰链别进灰尘液体}：虽然多数标称IPX8防水，但防尘等级普遍较低。沙滩、工地等环境要小心。
  \item \textbf{开合力度均匀}：别单手暴力甩开，双手从中间匀速展开最稳。
\end{itemize}

\textbf{内屏损坏的高发场景}\par

\begin{itemize}
  \item 合上时夹住耳机线/充电线——直接压裂内屏。
  \item 放包里和钥匙混装——外屏没事，内屏被钥匙硌出裂痕。
  \item 冬天从冷处进暖房立刻开合——温差导致玻璃应力变化。
  \item 用了一年后觉得"折痕变深了"去撕原厂膜——撕膜会连内屏一起损坏。
\end{itemize}

\booksection{六、折叠屏的“坑”和心理预期}

买之前心里要有这几本账：

\begin{itemize}
  \item \textbf{折痕会一直在}：2026年的技术做不到"完全无痕"，只能做到"亮屏正面几乎看不出"。如果你是完美主义者，每天盯着折痕难受，别买。
  \item \textbf{更重更厚}：横向折叠屏展开后普遍230g+，单手握持比直板机累；合上后厚度12-14mm，裤兜里更鼓。
  \item \textbf{续航普遍一般}：因为要塞铰链机构，电池空间被压缩，同价位续航通常不如直板机。
  \item \textbf{保值率低}：一年后二手价通常是新机的50-60\%，比直板机跌得快。
  \item \textbf{维修贵}：内屏总成维修价2000-4000元，铰链维修500-1500元。买之前查清楚保修政策。
  \item \textbf{App适配参差}：主流App已适配大屏，但小众App可能只是"拉伸放大"，体验一般。
\end{itemize}

折叠屏不是“更好的手机”，而是“另一种形态的手机”。买它是因为你需要它独特的形态，不是因为它“更先进”。

—— 关于折叠屏的选购哲学

\booksection{七、主流品牌怎么选（2026年）}

不写具体型号（会过时），只写品牌调性，帮你快速定位：

\begin{itemize}
  \item \textbf{华为Mate X系列}：横向折叠标杆，鸿蒙生态协同强，外屏比例接近正常手机，合上能当主力机用。价格高但保值率相对最好。
  \item \textbf{三星Galaxy Z Fold/Flip}：折叠屏鼻祖，铰链调校成熟，Z Flip是纵向折叠的颜值担当。国内售后网点较少是短板。
  \item \textbf{小米MIX Fold/Flip}：性价比路线，配置堆料足，价格相对亲民。适合想尝鲜又不想花太多的用户。
  \item \textbf{vivo X Fold系列}：横向折叠里外屏比例做得最接近正常手机的，合上当直板机用体验好。
  \item \textbf{OPPO Find N系列}：横向折叠里尺寸偏紧凑的，单手可握持，适合觉得大折叠太重的用户。
  \item \textbf{华为Pocket系列}：纵向折叠里外屏最大最实用的，化妆镜、防晒检测等女性向功能做得细致。
\end{itemize}

选购建议：预算充足追求生态选华为Mate X；想要最成熟铰链选三星Z Fold；性价比首选小米/vivo；纵向折叠看华为Pocket和三星Z Flip。

\begin{tcbinfo}{最后提醒}
折叠屏技术每年都在迭代，本附录写于2026年。买之前务必查最新评测和用户反馈，重点关注：铰链寿命实测、折痕表现、内屏维修政策这三项。线下亲手摸永远比看评测重要。

\end{tcbinfo}
